\documentclass[11pt]{article}
\usepackage{icmlko}
\begin{document}

\title{환율이 손잡이마다 다르다 --- 용량은 한 방향이 아니었다}
\author{월드모델 실험실}
\date{노트 459 · 2026년 8월 2일}
\maketitle

\begin{abstract}
노트 457 이 깊이를 $6$ 에서 $12$ 로 바꾸면서 손잡이 하나가 짝 없이
남았다 --- \textbf{lr $0.10$ 은 깊이 $6$ 에서 고른 값}이다(노트 427 ·
432). 깊은 나무는 보통 낮은 lr 을 원하므로 $0.10$ 아래를 촘촘히 봤다.

\textbf{안 움직인다.} 깊이 $12$ 에서도 lr $0.10$ 이다.

\begin{center}
\begin{tabular}{lrrrr}
\hline
lr & 판 & 날짜 통제 & KR 만화 & 비게임 앱 \\
\hline
$0.04$ & $+0.4759$ & $+0.5161$ & $+0.5578$ & $+0.2670$ \\
$0.06$ & $+0.4814$ & $+0.5209$ & $+0.5578$ & $+0.2667$ \\
$\mathbf{0.10}$ & $\mathbf{+0.4839}$ & $+0.5228$ & $+0.5561$ & $+0.2666$ \\
$0.16$ & $+0.4839$ & $+0.5234$ & $+0.5499$ & $+0.2671$ \\
\hline
\end{tabular}
\end{center}

씨앗 열을 짝지어 확인했다 --- lr $0.10 \to 0.06$ 은 판 $-0.0025$
$[-0.0038,-0.0013]$ 로 \textbf{확실히 잃는데} KR $+0.0019$
$[-0.0020,+0.0058]$ 와 앱 $+0.0012$ $[-0.0014,+0.0038]$ 은 \textbf{둘 다
$0$ 을 문다.} 노트 426 거울의 조건 ②가 막는다. \textbf{lr $0.10$ 확정.}

\textbf{그런데 표가 다른 것을 보여 준다.} lr 도 깊이와 \emph{같은
방향}으로 거래한다 --- 낮추면 판을 잃고 집 밖을 산다. 다만
\textbf{환율이 다르다.}
\end{abstract}

\begin{figure}[t]
\centering
\includegraphics[width=\textwidth]{figs/f1.pdf}
\caption{\textbf{가} --- 판에서 낸 것 대 집 밖에서 산 것. 기울기가
환율이다. \textbf{나} --- lr 은 판만 확실히 잃는다.}
\end{figure}

\section{환율표}

노트 458 이 규약 구멍 하나를 적어 뒀다 --- \emph{노트 426 거울에 거래
한도가 없다.} 판 $-0.02$ 에 집 밖 $+0.04$ 면 넣나? 그 물음에 답하려면
먼저 \textbf{지금까지의 거래가 어떤 값에 일어났는지}를 알아야 한다.
씨앗 열을 짝지어 잰 것만 모았다.

\begin{center}
\begin{tabular}{lrrrl}
\hline
변경 & 판 차 & 집 밖 평균 & \textbf{환율} & 판정 \\
\hline
\texttt{OOC\_DROP} (노트 441) & $0.0000$ & $+0.3681$ & $\infty$ & 채택 \\
\textbf{깊이 $6 \to 12$} (노트 457) & $-0.0032$ & $+0.0063$ & $\mathbf{1.95}$ & 채택 \\
lr $0.10 \to 0.06$ & $-0.0025$ & $+0.0016$ & $0.62$ & 기각 \\
깊이 $12 \to 16$ (노트 458) & $+0.0000$ & $+0.0001$ & --- & 기각 \\
K $32 \to 64$ (노트 456) & $+0.0002$ & $+0.0001$ & --- & 기각 \\
\hline
\end{tabular}
\end{center}

\textbf{깊이는 lr 보다 세 배 좋은 환율로 판을 집 밖과 바꾼다.} 그리고
\texttt{OOC\_DROP} 은 아예 공짜였다 --- 판을 $0.0000$ 내고
$+0.37$ 을 샀다.

\section{그래서 규약 구멍이 생각보다 작다}

노트 458 은 \emph{비율 한도}가 없어서 큰 거래를 못 막는다고 걱정했다.
표를 보면 걱정의 방향이 조금 달랐다. \textbf{조건 ②(집 밖 구간이 $0$ 을
안 문다)가 이미 대부분을 막는다} --- lr $0.06$ 은 환율 $0.62$ 로 나쁜
거래인데, 규약은 \emph{환율}이 아니라 \emph{구간}으로 그것을 걸렀다.
값이 작은 거래는 대개 구간이 $0$ 을 물기 때문이다.

그래서 한도를 새로 만드는 대신 \textbf{기록 의무}를 붙인다.

\begin{quote}
\textbf{조항} --- 판을 잃고 집 밖을 사는 변경은 \textbf{환율}(집 밖
평균 이득 $\div$ 판 손해)을 같이 적는다. 규약이 그 값으로 막지는 않지만,
\textbf{나중 거래는 앞선 거래의 환율과 견줄 수 있어야 한다.} 지금까지의
가장 나쁜 채택 환율은 $\mathbf{1.95}$(깊이 $12$)다.
\end{quote}

\section{용량은 한 방향이 아니다}

노트 $408 \sim 410$ 은 깊이 · lr · 자루섞기 셋이 \textbf{같은 용량
가족}이라고 적었다 --- 도메인이 그 셋에 매기는 선호 순서가 같고,
정체는 ``F18$-$선형 격차''(rho $+0.773$)라는 것이다.

이번 결과가 그 문장을 좁힌다. \textbf{선호 순서는 같아도 최적점은 따로
움직이지 않는다.} 깊이를 두 배로 늘려도 lr 의 봉우리는 $0.10$ 에
그대로 있다. 셋이 완전한 대체재라면 깊이를 올렸을 때 lr 이 내려가야
하는데 안 내려간다.

그리고 \textbf{환율이 셋 다 다르다.} 같은 ``용량''이라도 깊이로 산
집 밖 이득과 lr 로 산 것의 값이 세 배 차이다. \emph{용량은 크기 하나로
적을 수 있는 것이 아니라 방향이 있는 것}이다.

\section*{남은 것}

\begin{itemize}
\item \textbf{it(반복 수)은 아직 안 봤다} --- $220$ 은 언제 고른 값인지
기록이 없다. 깊이 $12 \cdot$ lr $0.10$ 에서 다시 볼 손잡이가 하나 남았다.
\item \textbf{환율이 왜 다른가.} 깊이는 \emph{상호작용}을 늘리고 lr 은
\emph{보폭}을 줄인다. 집 밖에서 값이 있는 것이 상호작용이라면 그게 답이다
--- 아직 안 쟀다.
\item \textbf{집 밖 짝 셋째}(노트 458) --- 둘뿐이면 조건 ①이 너무 쉽다.
\end{itemize}

\end{document}
